\chapter{اصول راهنمای طراحی}
پیش از ورود به جزئیات ساختارهای پیشنهادی، لازم است اصول بنیادینی که طراحی کل نظام نهادی گذار بر مبنای آن‌ها صورت گرفته، به‌روشنی بیان شود. این اصول برخاسته از ادبیات نظری سیاست تأسیسی (\lr{Constituent Politics})، تجربه‌ی تطبیقی گذارهای دموکراتیک و ضرورت‌های ویژه‌ی ایران است.

\begin{principlebox}{ترکیب دموکراسی مستقیم و تخصص‌گرایی (\lr{تکنوکراسی})}
هر تصمیم سیاسی دو بُعد دارد: بُعد واقعیت (\lr{Fact}) و بُعد ارزش (\lr{Value}). هنر حکمرانی خوب در دوران گذار، جدا کردن و سپس ترکیب درست این دو بُعد است.
\end{principlebox}

\begin{itemize}
    \item \rl{\textbf{بُعد واقعیت و تخصص \lr{(What is)}}: داده‌ها، حقایق علمی-تجربی، تحلیل‌های فنی، پیامدهای اقتصادی، حقوقی و اجتماعی هر تصمیم. این بُعد از تکنوکرات‌ها، متخصصان، مراکز پژوهشی و مشاوران دریافت می‌شود.}
    \item \rl{\textbf{بُعد ارزش و خواست ملی \lr{(What ought to be)}}: تمایلات، ارزش‌ها، آرزوها، ترجیحات و حق تعیین سرنوشت مردم. این بُعد \textbf{فقط و فقط} از طریق مراجعه‌ی مستقیم به آرای مردم (رفراندوم، نظرسنجی، مشاوره‌ی عمومی) قابل تعیین است.}
\end{itemize}

\textbf{مبنای نظری:} این تفکیک ریشه در اندیشه‌ی دیوید هیوم (تفکیک «هست» از «باید»)، یورگن هابرماس (دموکراسی مشورتی و کنش ارتباطی)، و فیلیپ پتیت (جمهوری‌خواهی و آزادی به‌مثابه عدم سلطه) دارد.

\textbf{تجربه‌ی جهانی:}
\begin{itemize}
    \item \textbf{ایرلند (۲۰۱۶-۲۰۱۹):} مجمع شهروندان ترکیبی از شهروندان عادی و متخصصان بود. شهروندان داده‌ها را دریافت کرده و سپس بر اساس ارزش‌های خود تصمیم می‌گرفند.
    \item \textbf{فرانسه (۲۰۱۹-۲۰۲۰):} کنوانسیون شهروندی اقلیم متشکل از ۱۵۰ شهروند تصادفی.
\end{itemize}

\begin{principlebox}{نمایندگی نسبتی \lr{(Proportional Representation)}}
\rl{در جامعه‌ای به تنوع و تکثر ایران، هیچ صدایی نباید نشنیده بماند. سیستم انتخاباتی نسبی-فهرستی تضمین می‌کند که ترکیب مجلس مهستان آینه‌ی واقعی جامعه‌ی ایران باشد.}
\end{principlebox}

\begin{principlebox}{قانون اساسی‌گرایی دو مرحله‌ای}
\rl{خلأ قانونی خطرناک‌ترین تهدید دوران گذار است. مدل دو قانون اساسی (موقت + نهایی) تضمین می‌کند که از روز اول چارچوب حقوقی وجود داشته باشد.}
\end{principlebox}

\begin{principlebox}{اصلاح نه تخریب}
\rl{تجربه‌ی عراق نشان داد: انحلال کامل نهادها منجر به فاجعه می‌شود. ساختارهای اداری، نظامی و خدماتی کشور با رهبری جدید و اصلاح‌شده حفظ می‌شوند.}
\end{principlebox}

\begin{principlebox}{شفافیت و پاسخ‌گویی حداکثری}
\rl{تمام جلسات مجلس مهستان علنی خواهد بود. قدرت بدون نظارت، فاسد می‌شود.}
\end{principlebox}

\begin{principlebox}{فراگیری و تعهد به حقوق بنیادین}
\rl{ایران برای همه‌ی ایرانیان. هیچ گروه، قوم، جنسیت، مذهب یا ایدئولوژی از فرایند گذار حذف نخواهد شد.}
\end{principlebox}
